\begin{abstract}
Resistive random access memory (RRAM) is a promising emerging nonvolatile memory (NVM) due to its large resistance ratio in different switching states. To improve memory density and reduce cost-per-bit, multi-level cell (MLC) RRAM stores multiple bits in a single cell, compared to a single-level cell (SLC). However, random mismatch, process variation, and resistance shift lead to reliability issues, degrade the probability of correct read, and increase the bit error rate (BER). 
This paper presents a time-based sensing scheme for robust read operation and extends to multi-level sensing for SLC and MLC RRAM arrays. Bit line (BL) voltage is converted into time delay by a voltage-to-time converter (VTC) and compared with the implicit timing reference generated by a delay line. By detecting different states in the time domain, the proposed time-mode sense amplifier (TSA) requires no analog reference voltage or current, which is used in the conventional voltage-mode sense amplifiers (VSA) or current-mode sense amplifiers (CSA). Power gating is employed to enable the time sampling only at the sensing points to suppress the short-circuit current. A charge sharing-induced error compensation (CSEC) circuit is used to eliminate the charge sharing-induced voltage drop and expand the sense margin by 1.56$\times$.
Monte Carlo simulations in 40nm technology show that the proposed TSA improves read reliability and reduces BER by 3-4 orders of magnitude compared to conventional VSA and CSA. 
The proposed time-based sensing scheme operates from 0.7-1.2 V supply and consumes 49 fJ/bit for read operation under a nominal 1.2 V supply.

% \textcolor{red}{1.2 V}

\end{abstract}